
\documentclass[10pt,a4paper]{article}

% ---------- Packages ----------
\usepackage[margin=1.8cm,headheight=14pt]{geometry}
\usepackage[T1]{fontenc}
\usepackage{lmodern}
\usepackage{microtype}
\usepackage{enumitem}
\usepackage{booktabs}
\usepackage{tabularx}
\usepackage{xcolor}
\usepackage{hyperref}
\usepackage{fancyhdr}

% ---------- Colours and hyperlinks ----------
\definecolor{headingblue}{RGB}{30,65,105}
\definecolor{lightblue}{RGB}{235,242,249}

\hypersetup{
    colorlinks=true,
    linkcolor=headingblue,
    urlcolor=headingblue
}

% ---------- Page layout ----------
\setlength{\parindent}{0pt}
\setlength{\parskip}{5pt}
\renewcommand{\arraystretch}{1.25}

\setlist[itemize]{
    leftmargin=1.5em,
    itemsep=3pt,
    topsep=3pt
}
\setlist[enumerate]{
    leftmargin=1.7em,
    itemsep=5pt,
    topsep=4pt
}

\pagestyle{fancy}
\fancyhf{}
\fancyhead[L]{\small Team 40}
\fancyhead[R]{\small Checkpoint One Status Report}
\fancyfoot[C]{\thepage}
\renewcommand{\headrulewidth}{0.4pt}

% ---------- Section formatting ----------
\usepackage{titlesec}

\titleformat{\section}
  {\large\bfseries\color{headingblue}}
  {\thesection.}{0.5em}{}

\titleformat{\subsection}
  {\normalsize\bfseries\color{headingblue}}
  {\thesubsection}{0.5em}{}

\titlespacing*{\section}{0pt}{12pt}{5pt}
\titlespacing*{\subsection}{0pt}{8pt}{3pt}

\begin{document}

% =========================================================
% TITLE
% =========================================================

\begin{center}
    {\LARGE\bfseries Checkpoint One Status Report}\\[5pt]
    {\large A Comprehensive Analysis of the Indian Startup Ecosystem}\\[8pt]
    {\normalsize\textbf{Team 40}}
\end{center}

\vspace{3pt}
\hrule
\vspace{5pt}

% =========================================================
% PAGE 1: METADATA AND PROGRESS
% =========================================================

\section{Meta-Information}

\begin{tabularx}{\textwidth}{@{}>{\bfseries}p{3.7cm}X@{}}
Open Project Title &
A Comprehensive Analysis of the Indian Startup Ecosystem \\

Group Number & 40 \\

Mentor Name & Duggineni Venkata Yugesh \\

Team Members &
Avanith (25B1013), Shanmukha (25B0907),
Tejesh (25B0905), Anirudh (25B1021) \\

GitHub Repository &
\textit{[Insert repository link]} \\
\end{tabularx}

\section{Progress}

\subsection{Datasets}

\textbf{Dataset 1: Startup Funding (2015--2021).}
Records of Indian startup investments containing investment dates,
startup names, industries, cities, funding amounts, and investment
stages.

\begin{center}
\small
\rowcolors{2}{lightblue}{white}
\begin{tabular}{c|rrrrrrr|r}
\toprule
\textbf{Year}
& 2015 & 2016 & 2017 & 2018 & 2019 & 2020 & 2021
& \textbf{Total}\\
\midrule
Records
& 4762 & 1018 & 685 & 310 & 111 & 978 & 1209
& 9073\\
\bottomrule
\end{tabular}
\end{center}

\textbf{Dataset 2: Indian Unicorns (May 2025).}
A curated list of 118 Indian unicorn startups, including location,
sector, time-to-unicorn, valuation, investors, and annual revenue.

\subsection{Data Cleaning and Preprocessing}

The following steps have been completed:

\begin{enumerate}
    \item Yearly funding files were standardised to a common
    seven-column schema and consolidated into \texttt{merged.csv}.

    \item Startup name variants were normalised. City names such as
    Bangalore and Bengaluru were consolidated, while cities outside
    India were grouped under ``Abroad''.

    \item Empty and irrelevant rows were removed. Investment stages
    were mapped to ten categories: Pre-Seed/Angel, Seed Funding,
    Pre-Series A, Series A, Series B, Series C, Series D \& Above,
    Private Equity, Debt/Bridge, and Undisclosed/Other.

    \item Funding amounts were converted to numeric values.
    Valuations, annual revenues, and total funding in the unicorn
    dataset were converted to crores of Indian rupees.
\end{enumerate}

\textbf{Updated Spreadsheet:}
\textit{[Insert spreadsheet link]}

\subsection{Division of Responsibilities}

Data collection and cleaning were performed collaboratively.
Avanith, Shanmukha, and Tejesh primarily handled Dataset 1,
while Anirudh focused on the analysis of Dataset 2.

\subsection{Challenges Encountered}

\begin{enumerate}
    \item \textbf{Uneven yearly coverage:} The dataset contains
    4,762 records for 2015 but only 111 for 2019, potentially
    distorting comparisons based on raw deal counts.

    \item \textbf{Inconsistent industry labels:} The 2015
    Industry/Vertical entries are more elaborate than those in
    subsequent years, complicating sector-wise standardisation.

    \item \textbf{Incomplete dates:} Where the exact day was
    unavailable, the first day of the corresponding month was
    assigned as a placeholder. These dates are approximate and
    unsuitable for precise daily-level analysis.
\end{enumerate}

% =========================================================
% PAGE 2: RESEARCH AND ANALYSIS
% =========================================================

\newpage

\section{Refined Problem Statement}

This study examines the geographical and sectoral distribution
of startup investment in India between 2015 and 2021, alongside
the characteristics and funding trajectories of startups that
eventually attained unicorn status.

\subsection{Research Question 1: Evolution of Startup Funding}

\textit{How did the distribution of startup investment across
Indian cities and sectors evolve between 2015 and 2021?}

The analysis will compare cities' shares of disclosed funding
and deal counts to assess changes in Bangalore's relative position
alongside Delhi-NCR, Mumbai, and emerging hubs. It will also
identify leading recipient sectors and examine shifts in their
funding shares, particularly around the COVID-19 period.

\subsection{Research Question 2: Characteristics of Unicorns}

\textit{What funding and business characteristics are associated
with Indian startups that achieve unicorn status?}

The analysis will investigate the time required to reach a
US\$1 billion valuation, the geographical and sectoral distribution
of unicorns, and variation in their reported valuations and
annual revenues.

Investment stages will also be examined to understand funding
trajectories as startups progress towards unicorn status. Comparing
early-stage and late-stage investments may reveal common financing
patterns and differences across sectors and cities.

\section{Methodology}

Both datasets will be analysed using descriptive statistics,
comparative analysis, and visualisation.

\begin{itemize}
    \item \textbf{Funding analysis:} Use total disclosed funding
    and deal counts to compare cities, sectors, and years.
    Transactions with undisclosed amounts will count as deals
    but will be excluded from monetary totals.

    \item \textbf{Unicorn analysis:} Examine time-to-unicorn,
    location, sector, valuations, revenues, and investment-stage
    patterns.

    \item \textbf{Temporal analysis:} Investigate changes in
    funding activity and composition around the COVID-19 period.
\end{itemize}

The study is exploratory: it aims to identify patterns and
associations rather than establish causal relationships.

\section{Limitations}

\begin{itemize}
    \item Funding coverage varies substantially across years,
    and many transactions have undisclosed amounts.

    \item Inconsistent labels and approximate investment dates
    may affect comparisons and temporal precision.

    \item The unicorn dataset is subject to survivorship bias
    because it includes only successful companies. Its financial
    figures may also reflect different reporting dates.
\end{itemize}

\section{Next Steps}

The next phase will involve city-wise, sector-wise, and year-wise
funding analysis; investigation of changes around COVID-19;
analysis of unicorn characteristics and investment stages; and
the preparation of visualisations with appropriate consideration
of data quality and coverage.

\end{document}
